\chapter{Implementation and Linear Solves}

\texttt{main.cpp} validates run parameters. \texttt{WaveEquation} keeps the
laboratory sequence: setup, matrix assembly, initialization, time loop,
diagnostics, and output. The PDE assembly, manufactured data, and Newmark
predictor/corrector are shared by every solver option.

\section{Why CG Is Applicable}
For a nonzero coefficient vector on the unconstrained space,
\[
 z^TMz=\|z_h\|_{L^2}^2>0,\qquad
 z^TKz=\|\nabla z_h\|_{L^2}^2\geq0.
\]
Both matrices are symmetric. On \(V_h\subset H^1_0(\Omega)\), the second
quadratic form is also strictly positive: a constant function with zero
boundary trace must be zero. Therefore
\[
 z^TBz=\|z_h\|_{L^2}^2+\tfrac14\Delta t^2\|\nabla z_h\|_{L^2}^2>0.
\]
\texttt{apply\_boundary\_values(...,true)} eliminates both rows and columns,
retaining positive diagonal entries on boundary DoFs. The solved matrices are
an SPD interior block plus a positive diagonal boundary block. Merely
eliminating rows would not justify the same symmetry argument.

Setting \texttt{WAVE\_CHECK\_MATRIX=1} checks the sparse entries for symmetry,
positive diagonals, and eight deterministic random Rayleigh quotients for
each constrained matrix. These checks can detect mistakes; they are not a
proof of positive definiteness. No dense conversion or eigendecomposition
is used.

\section{Solvers and Preconditioners}
The default is CG with SSOR at relaxation one. SSOR applies forward and
backward triangular sweeps to reduce the conditioning difficulty seen by CG.
For an SPD matrix and \(0<\omega_{\rm SSOR}<2\), the symmetric SSOR
preconditioner is SPD, as required by preconditioned CG.
Jacobi uses the inverse positive diagonal and is also SPD.
At relaxation one, writing \(B=D+L+L^T\), the SSOR matrix is
\(P=(D+L)D^{-1}(D+L)^T\), which makes its positive definiteness explicit.

The executable also accepts GMRES with right preconditioning, using a restart
basis of 30 in deal.II 9.5.1. Right preconditioning makes its stopping
criterion use the original residual, matching CG. All runs start each
linear solve from zero and use at most 5000 iterations, absolute tolerance
\(10^{-14}\), or relative reduction \(10^{-10}\).
After solving, the code independently computes the Euclidean residual
\(\|Ba-b\|_2\) (or \(\|Ma-b\|_2\)) and rejects an unsuccessful solve.
This algebraic residual is distinct from the spatial FEM error norms.

SSOR accounts for off-diagonal coupling, but costs triangular sweeps, depends
on ordering, is not mesh-independent, and is awkward to parallelize.
Jacobi is cheaper but ignores coupling. Neither is universally best.
Chapter 6 reports the measured comparison, including an important limitation:
a single smooth sine mode on a uniform mesh is unusually easy for Jacobi.

\section{Files and Reproduction}
The output files are:
\begin{itemize}
 \item \path{parameters.csv}: run parameters;
 \item \path{diagnostics.csv}: time, displacement errors and energy;
 \item \path{solver_stats.csv}: iterations, residuals and timings;
 \item \path{matrix_checks.csv}: sparse check data, only when enabled;
 \item \path{probe.csv}: centre displacement samples for phase/dispersion checks;
 \item \path{mesh.vtk}: mesh; PVD/VTU files: solution time series.
\end{itemize}
Python creates the plots.
\begin{verbatim}
bash scripts/docker_build_run.sh
source scripts/common.sh
load_course_modules
python scripts/review_validation.py
\end{verbatim}
The final two optional executable arguments select \texttt{cg|gmres} and
\texttt{ssor|jacobi}. Old argument lists retain the default.
Each run replaces its named CSV/PVD and matching snapshots.
Other files, including older unreferenced snapshots, are retained; open the
PVD rather than collecting VTU files by wildcard.
